
\section{Discussion}
\label{sec:discussion}

\subsection{Missing evidence is not the same as unknowably missing evidence}

The contrast between F1--F5 and F6 shows why missing evidence is too coarse a safety
category. The same formal condition---an obligation is uncovered---has different downstream
consequences depending on whether the remaining surface advertises the gap.

At minimum, risk depends jointly on (i) which obligation is missing, (ii) whether the
remaining surface exposes that omission, (iii) how the downstream policy responds to
visible gaps, and (iv) whether a separate coverage signal exists. This also explains why
aggregate mandatory-fact recall is not strictly monotonic with UDR across C2--C4.

\subsection{Superseding evidence creates an asymmetric omission}

Override and invalidation records are asymmetric: a later record can cancel the
actionability of an earlier one. If the later record is omitted, the older record does not
necessarily look incomplete. It can look \emph{more} coherent because the contradiction
has disappeared. This differs from many ordinary missing-fact cases, where the remaining
surface can still expose a hole.

This motivates treating superseding evidence as a first-class safety obligation in
long-lived decision systems. The relevant question is not only whether an old record is
stale, but whether the system can establish that all state-changing records required to
authorize a current decision have been surfaced.

\subsection{Relation to stale-memory and revocation-enforcement failures}

The nearest prior work sharpens rather than erases our boundary. STALE shows that an agent
may have access to later evidence yet fail to infer that an older belief is no longer valid
\citep{chao2026stale}. Revocation-enforcement work shows that a memory backend may return a
record even after it has been explicitly marked revoked \citep{shen2026revoked}. Both are
important, but both retain some representation of the update or revocation in the relevant
state or retrieval path.

Our controlled F6 intervention removes the later override/invalidation from the
decision-visible surface. The older instruction then need not be contradictory,
stale-marked, or revoked-marked from the policy's perspective. The failure is one of
\emph{missing transition evidence}, not merely incorrect adjudication of visible
transition evidence.

This also changes the mitigation boundary. A better state resolver helps when the update is
visible. A revocation filter helps when an invalid record is explicitly marked as such.
Obligation coverage addresses the complementary question: has the evidence acquisition
layer surfaced the classes of evidence required to justify a consequential decision?

\subsection{Mechanism identification, not prevalence estimation}

The benchmark deliberately constructs controlled worlds and a deterministic policy so
evidence visibility can be intervened on without stochastic model confounds. This gives
strong internal control but narrow external scope. The +16.67 percentage-point aggregate
effect is not a claim that 16.67\% of real agent decisions fail this way, and the 100\% F6
direct-treatment rate is not a field prevalence estimate.

The contribution is a reproducible mechanism demonstration and boundary characterization:
under a fixed policy with explicit visible-gap checks, some obligation omissions are
self-revealing and others are not; top-$k$ retrieval can turn additional families into
silent surfaces; and independently supplied completeness information changes behavior on
the same retrieved content.

\subsection{Why obligation coverage differs from policy confidence}

A policy-confidence score asks how confident the policy is \emph{given the evidence it
received}. Silent omission is precisely a case in which the received evidence can be
locally coherent enough to support high confidence.

An obligation-coverage interface asks a different question: whether the evidence
acquisition layer satisfied a predeclared set of evidence requirements. A practical
architecture may therefore need to separate:
\begin{enumerate}
  \item retrieval relevance---which records appear useful;
  \item obligation coverage---whether required evidence classes are represented; and
  \item decision policy---what action is justified given the covered evidence.
\end{enumerate}

\subsection{C5 moves rather than solves the hard problem}

The central limitation of C5 is also the next research problem. In the benchmark,
\texttt{coverage\_complete} is computed from a frozen formal world specification. Real
systems rarely possess an oracle inventory of every fact that ought to exist.

The next technical problem is therefore not simply to add a boolean. It is how to
construct trustworthy evidence obligations and coverage monitors from schemas, workflows,
policies, provenance graphs, temporal update channels, or learned detectors without
reintroducing silent failure at the monitoring layer. A coverage monitor can itself fail
silently; false-complete and false-incomplete signals must be studied before any deployment
claim.

\subsection{Why P3 is not required for the present claim}

The confirmatory claim concerns a compositional property of a frozen evidence surface and a
frozen downstream policy. P2 identifies that mechanism through controlled evidence
interventions with no stochastic model confound. A model-based replication could test
whether similar failure modes occur in independently chosen LLM policies, but it is not
required to establish the narrower deterministic causal result reported here.

If a later P3 replication is run, model identities and prompts must be frozen before
outputs and the results must remain a replication layer rather than a retroactive
modification of P2.
